% r03_fill_K4: proof of (K_0), part VI: the source-linear constant vanishes.

\subsubsection*{The source-linear constant vanishes}

Throughout, unit scale of level $r$ (\Cref{rot:def:eff,rot:def:match}): fine lattices $\Lambda_X=\eta'\Z^3$, $\Lambda_Y=R^{-1}(\eta'\Z^3)$, $\eta'=5^{-r}$, $\eta=5\eta'$; $\tau$ the cell-wise bijection with $|\tau x'-x'|\le c_\tau\eta'$; for $F\in\{X,Y\}$ the full propagator $C^F$ (mass $m_r=m\delta_r$; $C^Y=C^X\circ R$), the photon propagator $D^F_C$ (three scalar copies, frame-free), the $r$-step fluctuation kernels $G^F=\sum_{j<r}S^F_j$ and $D^F$, and the coarse covariances $\Gamma^F=C^F-G^F=H^FB^FH^{F*}$, $\Gamma^F_{\rm ph}=D^F_C-D^F$ (\Cref{rot:lem:Gam}(i)). Running couplings $\lambda_r=\lambda\delta_r$, $e_r^2=e^2\delta_r$; the $h$-vertex at zero coarse fields is $\delta_r^2\sum_x\eta'^3h(x)[|\phi'(x)|^2-C^F(x,x)]$, $N=2$ real components, with the free-boundary subtraction $C^F(x,x)=C(0)$ common to both flows. A unit-scale density $d$ corresponds to the physical density $\delta_r^{-1}d$. Constants depend on $a,m,\lmu,N,\tau$ only. $L_\eta\defeq C\log^3(1/\eta)$.

\emph{Kernel toolbox} ($u$ the displacement of the two arguments; ``both flows'' $=F\in\{X,Y\}$).
(Q1) $|\nabla^jC(u)|\le C_j(|u|+\eta')^{-1-j}e^{-cm_r|u|}$, $j\le3$, lattice differences along frame directions; the same for $D_C$; constants uniform in $\eta'\le1$ and in the mass ($m_r\in[0,m]$ for $C$, $\lmu\delta_r\in[0,\lmu]$ for $D_C$; $c$ halved for $j=3$). This is the standard bound for the massive free lattice Laplacian: scale to $\eta'=1$, $\mu=m_r\eta'$, symbol $E(q)+\mu^2$ with $E(q)=\sum_\nu4\sin^2(q_\nu/2)$; by the lattice symmetries let $|u_1|=|u|_\infty$, and shift the $q_1$-contour by $\mathrm i\kappa$, $\kappa=\min(\mu/2,1)$, where $\Re[E(q+\mathrm i\kappa e_1)+\mu^2]\ge\frac12(E(q)+\mu^2)$ ($\cosh\kappa-1\le\kappa^2$; the vertical sides cancel by periodicity): factor $e^{-\kappa|u_1|}$; on the shifted contour the integrand with $j$ differences obeys $|\partial_{q_1}^b(\cdot)|\le C_b(E+\mu^2)^{(j-2-b)/2}\le C'_b|q|^{j-2-b}$ uniformly in the mass ($j\le2$), and a cutoff at $|q|\sim|u_1|^{-1}$ with $j+2$ integrations by parts in $q_1$ on the high part gives $|u_1|^{-1-j}$ ($j=3$: an extra factor $1+\mu|u_1|$, absorbed by halving $c$). Both factors come from one integral, so the product of the two bounds holds.
(Q2) Matched-pair differences of $C$-type kernels: \Cref{rot:cor:C} gives $|C^X(y',z')-C^Y(\tau y',\tau z')|\le C_6\eta'^{1/2}|u|^{-3/2}$ for $|u|\ge25\eta'$ (the same for $D_C$); for a kernel carrying $j\le2$ differences, interpolation of this bound against the per-flow bounds $|\nabla^{j+1}C|\le C(|u|+\eta')^{-2-j}$ (difference quotients at an optimised step) gives the relative gains $(\eta'/|u|)^{1/4}$ ($j=1$) and $(\eta'/|u|)^{1/8}$ ($j=2$) on the per-flow size $(|u|+\eta')^{-1-j}$; inside $|u|<25\eta'$ the difference is bounded by the sum of the two per-flow sizes.
(Q3) Fluctuation kernels per flow: $|S^F_j(x,y)|\le C\ell_j^{-1}e^{-\delta|x-y|/\ell_j}$, $|\nabla S^F_j|\le C\ell_j^{-2}e^{-\delta|x-y|/\ell_j}$, $|\nabla S^F_j(x,y)-\nabla S^F_j(x',y)|\le C_{\alpha,r}\ell_j^{-2-\alpha}|x-x'|^\alpha e^{-\delta\min(|x-y|,|x'-y|)/\ell_j}$, $\ell_j=5^j\eta'$; hybrid flow: \Cref{rot:prop:BD}(a) at $A=0$, with $C_{\alpha,r}\le C_\alpha\log^{3\alpha}(2+\ell_j/\eta')$ on slice $j$; standard flow: \cite[Lemmas~2.1, 2.2, 2.4]{Ba4} and \cite[Thm.~3.3]{K1} at $A=0$. Hence $|\nabla^jG^F(u)|\le C(|u|+\eta')^{-1-j}e^{-\delta'|u|}$, $j=0,1$.
(Q4) Coarse covariances per flow: $\sup|\Gamma^F|\le C_1$ (\Cref{rot:lem:Gam}(ii)); $|\Gamma^F(x,y)|,|\nabla_y\Gamma^F(x,y)|\le C(1+|x-y|)^{-1}e^{-cm_r|x-y|}$, uniformly in $r$ and $m_r\in[0,m]$: for $|x-y|\ge1$ the triangle inequality between (Q1) ($j=0,1$) and the fixed-rate decay $|\nabla_y^iG^F(x,y)|\le C_Ge^{-\delta'|x-y|}$, $i\le1$, which (Q3) gives by summing the slices ($\ell_j\le\frac15$), with $e^{-\delta'\rho}\le c_{\delta'}\rho^{-2}e^{-(\delta'/2m)m_r\rho}$; for $|x-y|<1$ it is $C_1$ and the Lipschitz bound with $C'_\nabla\defeq(2\pi)^{-3}C_BC_hC_{\nabla h}\int_{V_1}|p|^{-2}dp$ from \Cref{rot:lem:Gam}(ii) and the bounds $|\hat h_x|\le C_h$, $|\nabla_y\hat h_y|\le C_{\nabla h}$ of the proof of \Cref{rot:lem:CK1}. H\"older $|\Gamma^F(x,y)-\Gamma^F(x,w)|\le C_2|y-w|^{1/2}$ (\Cref{rot:lem:Gam}(iv)) and Lipschitz with constant $\sqrt3C'_\nabla$. The same for $\Gamma_{\rm ph}$ (mass $\lmu\delta_r$).
(Q5) Second differences, both flows: $|\nabla\nabla H^F(x,z)|\le L_\eta$ and $|\nabla_y\nabla_y\Gamma^F(x,y)|\le CL_\eta$ (\Cref{rot:lem:CK1}).
(Q6) Matched-pair differences without derivatives: $|\Gamma^X(x',y')-\Gamma^Y(\tau x',\tau y')|\le C_\Gamma\eta^{1/2}$ (\Cref{rot:lem:Gam}(iii)), and $\le C_6\eta'^{1/2}|u|^{-3/2}$ for $|u|\ge25\eta'$ (\Cref{rot:cor:C}, for $G^F_r$ and $\Gamma^F$ separately); $|G^X(x',y')-G^Y(\tau x',\tau y')|\le C_6\eta'^{1/2}|u|^{-3/2}$ ($|u|\ge25\eta'$) and $\le2C_0(|u|+\eta')^{-1}$ always; the same for the photon kernels.
(Q7) Vertex matching: \Cref{rot:lem:vert}(i)--(iii). (Q8) Summation by parts on the infinite lattice and the lattice Leibniz rule $\nabla_\mu(FG)(y)=(\nabla_\mu F)(y)G(y+\eta'e_\mu)+F(y)(\nabla_\mu G)(y)$.

\begin{lemma}[the block-average symbol]\label{rot:lem:CK0}
The $(r-1)$-step centred cube average is the composition of $5$-point averages along each rotated axis, whose symbols telescope, $\prod_j\sin(5^{j+1}t/2)/(5\sin(5^jt/2))=\sin(5^{r-1}t/2)/(5^{r-1}\sin(t/2))$; hence $u_{r-1}(q)=\prod_\alpha\sin(q_\alpha/2)/(\eta^{-1}\sin(\eta q_\alpha/2))$, $q_\alpha=q\cdot f_\alpha$, on $V_\eta$. Consequently $J_0\defeq\int_{V_\eta}|u_{r-1}(q)|\,dq\le(C\log(1/\eta))^3$ (the $L^1$ norm of the Dirichlet kernel per axis), $J_1\le C$, and $\sum_{l_1}|u_{r-1}(p+l_1)|\le C\log^3(1/\eta)$, $\sum_{l_1\ne0}|u_{r-1}(p+l_1)|/|p+l_1|\le C$, uniformly in $p\in V_1$.
\end{lemma}
\begin{proof}
The telescoping of the Dirichlet kernels along each axis; the $L^1$ norm of the Dirichlet kernel of order $n$ is $\asymp\log n$.
\end{proof}

\begin{lemma}[second differences of the coarse elements, both flows]\label{rot:lem:CK1}
$|\nabla\nabla H^F(x,z)|\le CL_\eta$ for every pair of frame directions, $F\in\{X,Y\}$; and, uniformly in $m_r\ge0$, $|\nabla_y\nabla_y\Gamma^F(x,y)|\le CL_\eta$ and $|\nabla_y\nabla_w\Gamma^F(y,w)|\le C$. The same for $\Gamma_{\rm ph}$.
\end{lemma}
\begin{proof}
By \Cref{rot:lem:H}, $H^F(x,z)=(2\pi)^{-3}\int_{V_\eta}e^{\mathrm iq\cdot(x-z)}u_{r-1}(q)M^F(q)A^F(q)\,dq$ with $M^F(q)=\sum_le^{\mathrm i(q+l)\cdot r_F}u_F(q+l)/\Delta'_F(q+l)$. A second lattice difference in the fine variable multiplies the alias-$l$ summand by $d_F(q+l)^2$, $|d_F(Q)|\le|Q|$; since $\Delta'_F(Q)\ge c|Q|^2$ on every alias cell and $\sum_lu_F(q+l)^2=1$ (\Cref{rot:prop:fc}), $|\sum_ld_F^2u_F/\Delta'_F|\le C$, and the integral against $|u_{r-1}|$ costs $J_0\le CL_\eta$ (\Cref{rot:lem:CK0}). For $\Gamma$ write $\Gamma(x,y)=(2\pi)^{-3}\int_{V_1}B(p)\hat h_x(p)\overline{\hat h_y(p)}\,dp$, $\hat h_y(p)=\sum_zH(y,z)e^{-\mathrm ipz}$, which by Poisson summation equals $A^F(p)\sum_{Q=p+l_1+l}e^{\mathrm iQ\cdot y}u_{r-1}(p+l_1)u_F(Q)/\Delta'_F(Q)$; a fine difference in $y$ multiplies the $Q$-summand by $d_F(Q)$. Hence $|\hat h_y|\le C$, $|\nabla_y\hat h_y|\le C[|p|A^F/\Delta'_F+a\sum_{l_1\ne0}|u_{r-1}(p+l_1)|/|p+l_1|]\le C$ and $|\nabla_y\nabla_y\hat h_y|\le Ca\sum_{l_1}|u_{r-1}(p+l_1)|\sum_l|u_F|\le CL_\eta$, uniformly in $p\in V_1$ and $m_r\ge0$ (\Cref{rot:lem:CK0}). With $B^F(p)\le C_B|p|^{-2}$ (\Cref{rot:lem:Gam}(ii)), integrable on $V_1$ in three dimensions, both bounds follow.
\end{proof}

\begin{lemma}[matched-pair comparison of gradient vectors]\label{rot:lem:CK2}
Let $g$ on $\Lambda_X$ and $h$ on $\Lambda_Y$ satisfy, on a region, $|g(x')-h(\tau x')|\le A$ at matched pairs, $|\nabla g|,|\nabla h|\le C_\nabla$, and the H\"older modulus $|\nabla g(y)-\nabla g(y+v)|\le C_H|v|^\alpha$ (likewise $h$), $0<\alpha\le1$. With $\nabla^Xg(y')\defeq\sum_\mu(\nabla_\mu g)(y')e_\mu$ and $\nabla^Yh(\tau y')\defeq\sum_\mu(\nabla^Y_\mu h)(\tau y')f_\mu$ the lattice gradient vectors, for every $\rho\in\eta'\N$
\[
|\nabla^Xg(y')-\nabla^Yh(\tau y')|\le3\big[2A/\rho+2C_H\rho^\alpha+4c_\tau C_\nabla\eta'/\rho\big].
\]
\textup{(a)} For $\Gamma$-type kernels, $A=C_\Gamma\eta^{1/2}$ (Q6), $C_H=CL_\eta$ with $\alpha=1$ (Q5): $\rho=\eta^{1/4}$ gives $|\nabla^X\Gamma^X(x',\cdot)(y')-\nabla^Y\Gamma^Y(\tau x',\cdot)(\tau y')|\le C\eta^{1/4}L_\eta=:A_1$. \textup{(b)} Decaying form: $A=C_6\eta'^{1/2}|u|^{-3/2}$ for $|u|\ge25\eta'$ gives $\le C\eta'^{1/4}|u|^{-3/4}L_\eta$. \textup{(c)} Fluctuation kernels slice by slice: on slice $j$, $A_j=C\eta'^{1/2}\ell_j^{-3/2}e^{-\delta|u|/\ell_j}$, $C_{H,j}=C_\alpha\log^{3\alpha}(2+\ell_j/\eta')\,\ell_j^{-2-\alpha}$, $C_{\nabla,j}=C\ell_j^{-2}$; with $\alpha=\frac12$ and $\rho=\ell_j(\eta'/\ell_j)^{1/3}$, $|\nabla^XS^X_j(x',y')-\nabla^YS^Y_j(\tau x',\tau y')|\le C_\vartheta\ell_j^{-2}(\eta'/\ell_j)^{\vartheta}e^{-\delta|u|/(2\ell_j)}$ for every $\vartheta<\frac16$ and every $u$, and, summing over $j<r$, $|\nabla^XG^X(x',y')-\nabla^YG^Y(\tau x',\tau y')|\le C\eta'^{1/6}\log(1/\eta)\,(|u|+\eta')^{-13/6}e^{-\delta'|u|}$.
\end{lemma}
\begin{proof}
Lattice Taylor expansion: for $v\in\Lambda_Y-\Lambda_Y$, $h(\tau y'+v)-h(\tau y')=\nabla^Yh(\tau y')\cdot v+\varrho$, $|\varrho|\le3C_H|v|^{1+\alpha}$ (sum the increments along a lattice path of at most $3|v|/\eta'$ bonds). Take $v_\mu\defeq\tau(y'+\rho e_\mu)-\tau y'$, $|v_\mu-\rho e_\mu|\le2c_\tau\eta'$; then $e_\mu\cdot\nabla^Yh(\tau y')=\rho^{-1}[h(\tau(y'+\rho e_\mu))-h(\tau y')]+O(C_H\rho^\alpha)+O(C_\nabla c_\tau\eta'/\rho)$, and likewise $e_\mu\cdot\nabla^Xg(y')=\rho^{-1}[g(y'+\rho e_\mu)-g(y')]+O(C_H\rho^\alpha)$. Subtract: the bracket differences are matched-pair differences, each $\le A$; sum over the three components. In (c) the two main terms balance at the stated $\rho$, giving the exponent $\alpha/(2(1+\alpha))=\frac16$, the $\eta'$ term is smaller, and the sum over $j$ is dominated by $\ell_j\sim|u|$.
\end{proof}
